\section{System Architecture and Command Semantics}
\label{sec:architecture}
\subsection{Ownership Boundaries and Data Flow}
The original networking simulator remains in its existing repository modules.
The power-plane extension is housed separately under
\code{research_power_plane}, with dedicated command adaptation, topology,
timeline, PSCAD mapping, automation, measurement, and dashboard modules.
This separation allows networking colleagues to evolve their algorithms while
the adapter handles changes in their output fields. It does not guarantee
compatibility with an arbitrary future schema: new aliases and semantic changes
must be tested at the adapter boundary.

Figure~\ref{fig:architecture} shows the implemented workflow. A controller
output batch or a graphical dashboard submission supplies topology, commands,
optional faults, and logical-to-PSCAD bindings. Canonical JSON remains the
exchange representation. The dashboard provides feeder selection, switch-state
controls, scheduled operations, fault forms, and mapping tables; advanced raw
JSON remains available. Committing an edit can automatically invoke the same
pipeline as a command-line replay. Selecting an element or moving the
scheduled-state preview is not a simulation request.

\begin{figure}[tb]
\centering
\includegraphics[width=\textwidth]{architecture.pdf}
\caption{Implemented batch workflow. Solid arrows indicate current processing
paths. Measurement display follows run completion; no feedback path from the
running EMT solver to the networking controller is implemented.}
\label{fig:architecture}
\end{figure}

The displayed network state and the PSCAD wiring are related but not
interchangeable. The present dashboard operates on the fixed six-bus electrical
case and its component map. A new graph edge without a corresponding wired
PSCAD branch does not create a new physical circuit. Extending the feeder
requires an electrical-model and mapping update.

\subsection{Command Contract and Graph Model}
Represent a normalized command as
\begin{equation}
 c_k = (\mathrm{id}_k,\,\tau_k,\,a_k,\,\mathcal{O}_k,\,\mathcal{M}_k),
\end{equation}
where $\tau_k$ is a simulation-time timestamp, $a_k$ is an action,
$\mathcal{O}_k$ contains target nodes, edges, or a requested path, and
$\mathcal{M}_k$ holds controller identity and metadata. Supported actions
include opening or closing a line, isolating or restoring a node, applying
a switch set, and requesting a reroute path. Commands are processed in a
stable chronological order, retaining input order for equal timestamps.

Let the physical connectivity model be the undirected graph
$\mathcal{G}=(\mathcal{V},\mathcal{E})$, with $x_e=1$ for a logically closed
branch and $x_e=0$ for an open branch. The currently conducting graph is
\begin{equation}
 \mathcal{G}_x = (\mathcal{V},\{e\in\mathcal{E}:x_e=1\}).
\end{equation}
For a source set $\mathcal{S}$, the software's energized-node estimate is the
union of connected components containing a source. This is a connectivity
estimate, not a voltage calculation. Likewise, the analysis field called
``island count'' counts connected components; it does not establish viable
grid-forming operation within a separated component.

With radiality preservation enabled and an initially acyclic graph, a proposed
closure $e=(u,v)$ is rejected if a path already connects $u$ and $v$ in
$\mathcal{G}_x$. Opening an edge cannot introduce a cycle. This local rule
preserves a forest during accepted edge changes, but does not require every
load to remain connected. It also does not enforce voltage limits, branch
ratings, source capacity, protection selectivity, or synchronization between
multiple sources.

An implementation detail limits the strength of the validation guarantee.
The current command handler records command-level warnings while processing
individual requested operations. A command may therefore be marked
unaccepted or partially accepted after some valid edge operations have already
been applied. Only accepted state-changing operations enter the switching
timeline; rejected operations and no-ops do not. This is
\emph{not an atomic, all-or-nothing transaction}. The loop-rejection experiment
tests a single rejected closure, not general rollback behavior. A production
interlock would require validation of the complete prospective transition
before committing it.

\subsection{Schedule Compilation and Execution}
Each logical line identifier maps to a PSCAD timed-logic component and its
associated three-phase breaker. The component's control convention reverses
the graph variable:
\begin{equation}
 b_e(t)=1-x_e(t),\qquad b_e=0\text{ closed},\quad b_e=1\text{ open}.
\label{eq:polarity}
\end{equation}
The adapter writes the initial state and operation count and times using
\code{INIT}, \code{NUMS}, \code{TO1}, and \code{TO2}. The current mapping
permits at most two transitions per breaker during one run, consistent with
the selected timed component \cite{pscadTimedBreaker}. An unchanged branch
receives an operation beyond the simulation window to accommodate the
component's lack of a zero-operation setting. The generated waveform inside
the experiment window remains constant.

Execution uses a project-local wrapper around PowerMCP, a Model Context
Protocol (MCP) server exposing PSCAD automation tools. A persistent localhost
service owns the application connection. The local wrapper focuses the intended
case and uses a blocking project run so that returning control means the run
has completed. The persistence avoids reopening the application for every
dashboard action; it does not represent an independently powered,
always-available field control network.

The client serializes project parameter writes with a run lock. To avoid the
PSCAD output-overwrite dialog, it removes only the current case's generated
output before building. Other case files and separately named archives are
preserved. Successful execution requires acceptable build messages, a newly
produced output fingerprint, and the expected measured channels. Tool errors,
missing channels, or stale results fail the request rather than being presented
as a successful new simulation.

These are engineering safeguards, not a formal proof of end-to-end correctness.
In particular, clearing the current raw output means that publication-grade
reproducibility needs a separate immutable archive for every run. Each new
submission starts another continuous EMT experiment from its specified initial
conditions. The current system does not alter a running simulation in response
to newly arriving controller messages.
